\section{Every original section value under the fixed-model native analytic charts}
\label{sec:projective-fixed-native-analytic-values}

\begin{theorem}[Whole original section values commute with the actual analytic chart comparison]
\label{thm:projective-fixed-native-analytic-values}
\lean{EnvelopeProperty.ProjectiveFixedNativeAnalyticValues.nativeChartIso_section, EnvelopeProperty.ProjectiveFixedNativeAnalyticValues.nativeChartIsoAt_section, EnvelopeProperty.ProjectiveFixedNativeAnalyticValues.nativeAffineChartHomeomorph_val}
\leanok
\uses{thm:projective-fixed-native-model-charts,thm:projective-fixed-native-analytic-charts}
\source{bj22-trivval:sec:tc}
The actual native chart scheme isomorphisms pull every original section
of the reconstructed chart algebra back to precisely its original
native section, through the canonical native open top-section
identification. Consequently the actual fixed-model native affine
analytic chart homeomorphism preserves the value of every entire
original section of the original native chart. This holds on the whole
original support-open valued-point domain and uses no section-value
matching premise.
\end{theorem}
\begin{proof}
For any actual ring equivalence, naturality of the original
global-section/spectrum isomorphism identifies the pullback of the
section of a mapped ring element with the section of its original
element. Apply this to the native section equivalence defining the
original chart scheme isomorphism. The canonical affine-open spectrum
isomorphism then pulls that section back to its exact original native
open section. Composing with the literal equality equivalence between
the local reconstruction chart and the same chart of the one fixed
model gives the identical statement for the fixed-model chart scheme
isomorphism.

Express an original native support-open point as the restriction image
of an actual native point of the original open subscheme. Its image
under the constructed affine analytic chart is the actual spectrum
point obtained by the proved native chart scheme isomorphism. The
canonical native affine spectrum evaluation, true native point-map
evaluation, and the preceding exact section pullback identity identify
its value with the original open section evaluation. Finally the actual
native open-inclusion evaluation/top-section identity identifies this
with the value of the same original section at the original ambient
native point. Every entire original section, rather than a selected
coordinate family, is covered by this identity.
\end{proof}
